\documentclass[runningheads]{llncs}
\usepackage[T1]{fontenc}
\usepackage{graphicx}
\usepackage{amsmath,amssymb}
\usepackage{booktabs}
\usepackage{subcaption}
\usepackage{color}
\renewcommand\UrlFont{\color{blue}\rmfamily}

\begin{document}

\title{Hybrid CNN-LSTM Model for Short-Term Solar Power Forecasting with Anomaly-Aware Preprocessing}
\titlerunning{Hybrid CNN-LSTM Solar Forecasting with Anomaly-Aware Preprocessing}

\author{Tushar Basak \and Varun Saxena}
\authorrunning{T. Basak and V. Saxena}

\institute{Department of Electronics \& Communication Engineering, Faculty of Technology, University of Delhi, Delhi 110007, India\\
\email{tusharbasak555@ece.du.ac.in, varunsaxena@fot.du.ac.in}}

\maketitle
\markboth{Hybrid CNN-LSTM Solar Forecasting with Anomaly-Aware Preprocessing}
{Hybrid CNN-LSTM Solar Forecasting with Anomaly-Aware Preprocessing}

\begin{abstract}
Accurate ultra-short-term photovoltaic (PV) power forecasting depends on both model architecture and telemetry quality. We investigate an anomaly-aware preprocessing pipeline for 30-minute-ahead PV power forecasting using a one-year, 5-minute-resolution NREL dataset from a 70~MW utility-scale PV site, synthetically injecting three telemetry corruption types -- daytime dropouts, spikes, and short-duration offset faults -- across $\sim$2\% of observations. An Isolation Forest (IF) detects anomalies, reconstructed via local linear interpolation. To separate preprocessing from detection quality, we evaluate four parallel conditions: a never-corrupted \emph{pristine} ceiling, \emph{raw corrupted} telemetry, \emph{IF}-cleaned telemetry (proposed), and \emph{oracle}-cleaned telemetry. The detector achieves 72.21\% recall and 57.48\% precision. All deep learning results (mean $\pm$ std.\ dev.\ over 5 seeds) are validated with paired $t$-tests, Holm--Bonferroni correction, and a Two One-Sided Tests (TOST) equivalence procedure. The proposed CNN--LSTM achieves RMSE $=3.57\pm0.04$~MW, MAE $=1.64\pm0.08$~MW, $R^2=0.949\pm0.001$, statistically indistinguishable from pristine ($3.58\pm0.06$~MW) and oracle-cleaned ($3.58\pm0.05$~MW) (Holm-adjusted $p=1.000$ both), confirmed by TOST equivalence within a $\pm0.10$~MW margin ($p\approx0.05$). Against raw-corrupted data ($3.89\pm0.11$~MW), cleaning yields an 8.2\% mean RMSE reduction. The proposed model beats a plain LSTM in every one of 5 paired seeds ($3.77\pm0.11$~MW; raw $p=0.018$, Holm-adjusted $p=0.053$). Both deep models outperform persistence (4.93~MW) and a diurnal-lag autoregressive baseline (4.83~MW) by a wide margin.

\keywords{Anomaly detection \and CNN-LSTM \and Isolation Forest \and photovoltaic forecasting \and equivalence testing.}
\end{abstract}

\section{Introduction}
Solar irradiance is intrinsically stochastic, so PV plant output fluctuates unpredictably, jeopardizing voltage stability, power reserves, and dispatch schedules. Deep learning, particularly LSTM networks \cite{ref1}, has largely replaced classical statistical methods for solar forecasting due to its ability to capture non-linear temporal dependencies \cite{ref2,ref3}; hybrid CNN-LSTM architectures further improve performance by pairing 1D-CNN layers (local ramp dynamics) with LSTM layers (diurnal memory) \cite{ref4,ref5,ref6,ref7}. Most published literature, however, assumes pristine, noise-free telemetry, whereas real PV plants regularly suffer \textbf{sensor dropouts}, \textbf{glitch spikes}, and \textbf{short-duration offset faults} that distort gradient updates and degrade generalization. While anomaly detection in smart-grid telemetry has been studied independently \cite{ref8}, few works formalize an end-to-end framework combining unsupervised mitigation with hybrid CNN-LSTM forecasting, isolate detection quality from imputation quality, or subject the resulting claims to formal multiple-comparison and equivalence testing.

\noindent\textbf{Contributions.} (1) \emph{End-to-End Fault-Tolerant Architecture}: a pipeline coupling unsupervised IF outlier filtering and temporal interpolation with a hybrid 1D CNN-LSTM. (2) \emph{Realistic Fault Injection Protocol}: a reproducible synthetic corruption protocol modeling three PV fault modes across $\sim$2\% of telemetry rows. (3) \emph{Decomposed Four-Condition Ablation}: pristine and oracle-labeled conditions evaluated alongside the proposed pipeline, separating detection quality from imputation quality and quantifying recoverable error. (4) \emph{Rigorous Multi-Seed Statistical Validation}: mean $\pm$ std.\ dev.\ over 5 seeds, paired $t$-tests, Holm--Bonferroni correction, and TOST equivalence -- rather than inferring equivalence from a non-significant difference alone.

\section{Related Work}
Early solar forecasting relied on ARIMA/exponential smoothing \cite{ref9}, which struggle with abrupt ramps, and ML regressors (SVR, Random Forests) \cite{ref10} that lack recurrent mechanisms. LSTM/GRU networks \cite{ref1} overcome the vanishing-gradient problem of vanilla RNNs, and hybrid CNN-LSTM models \cite{ref4,ref5,ref6,ref7,ref11} extend this using 1D convolutions ahead of a sequential LSTM stage, sometimes with attention \cite{ref6,ref11} or decomposition preprocessing \cite{ref7} -- directions not pursued here, which instead targets the upstream data-quality problem. Anomaly detection in power systems has developed through statistical checks, autoencoders, and tree-based ensembles \cite{ref8}. The Isolation Forest \cite{ref12} offers strong efficiency on high-dimensional time series without distributional assumptions and has been applied to PV fault detection \cite{ref13}, but is overwhelmingly treated as a standalone monitoring problem; few works integrate it directly into a deep hybrid forecaster's training pipeline or decompose the benefit against oracle-labeled and pristine reference conditions using equivalence testing rather than non-significant $p$-values.

\section{Methodology}
The pipeline comprises site selection, cyclical feature engineering, synthetic anomaly injection, parallel construction of four data conditions (pristine, raw-corrupted, IF-cleaned, oracle-cleaned), sliding-window sequencing, and CNN-LSTM modeling evaluated identically across all conditions, followed by a multi-seed ($n=5$) benchmark with paired $t$-tests, Holm--Bonferroni correction, and TOST equivalence testing.

\noindent\textbf{Dataset and Site Selection.} Data are drawn from the NREL Solar Power Data for Integration Studies (2006, SQLite, 274 PV stations across Alabama). A strict filter selects a single UPV facility to avoid cross-site contamination: latitude $30.45^\circ$N, longitude $-88.25^\circ$W, rated capacity 70.0~MW, $\Delta t=5$~min, $N=105{,}120$ observations (365$\times$288), zero missing intervals.

\noindent\textbf{Cyclical Feature Engineering.} Since temperature/GHI sensors are absent, four cyclical features are derived from timestamp $t$:
\begin{equation}
\text{hour\_sin}(t)=\sin\!\Big(\tfrac{2\pi h(t)}{24}\Big),\;
\text{hour\_cos}(t)=\cos\!\Big(\tfrac{2\pi h(t)}{24}\Big)
\end{equation}
\begin{equation}
\text{day\_sin}(t)=\sin\!\Big(\tfrac{2\pi d(t)}{365.25}\Big),\;
\text{day\_cos}(t)=\cos\!\Big(\tfrac{2\pi d(t)}{365.25}\Big)
\end{equation}
where $h(t)\in[0,24)$ is fractional hour-of-day and $d(t)\in[1,365]$ is day-of-year.

\noindent\textbf{Synthetic Anomaly Injection.} $\sim$2.0\% of observations ($N_{\text{anom}}=2{,}091$) are corrupted, in equal proportion: (1) \emph{Sensor Dropout} ($n_1=695$): forced to 0.0~MW during daylight ($P(t)>5.0$~MW); (2) \emph{Glitch Spikes} ($n_2=698$): scaled by $\alpha\sim\mathcal{U}(3.0,5.0)$, clipped to 1.5$\times$ capacity; (3) \emph{Short-Duration Offset Fault} ($n_3=698$): 4--12 consecutive intervals (20--60~min) receive additive offset $\delta\sim\mathcal{U}(0.15,0.35)\times$capacity. The pristine series is retained unmodified as the reference ceiling.

\noindent\textbf{Detection, Oracle Labeling, and Reconstruction.} An IF is trained on $[P_{\text{raw}}(t),\text{hour\_sin}(t),\text{hour\_cos}(t)]$; path length $h(x)$ over $T=100$ trees yields anomaly score
\begin{equation}
s(x,n)=2^{-\mathbb{E}(h(x))/c(n)}
\end{equation}
where $c(n)$ is the average unsuccessful-search path length of a Binary Search Tree. Points with $s(x,n)>\tau$ ($\gamma=0.025$) are masked and reconstructed by local linear interpolation:
\begin{equation}
\hat{P}_{\text{clean}}(t)=P(t_a)+\frac{P(t_b)-P(t_a)}{t_b-t_a}(t-t_a)
\end{equation}
bounded to $[0,P_{\max}]$. Against ground truth, the detector achieved \textbf{72.21\% recall} and \textbf{57.48\% precision} (2{,}627 points flagged); the moderate precision reflects a deliberately permissive $\gamma=0.025$ favoring recall over false positives near dawn/dusk ramps. The \textbf{oracle-cleaned} condition masks the exact injected indices with identical interpolation, isolating residual error from imperfect detection.

\noindent\textbf{Sliding Sequence Windowing.} Each condition is windowed with $W=24$ steps (2~h) and horizon $H=6$ (30~min):
\begin{equation}
\mathbf{X}_\tau=[\mathbf{x}_{\tau-W+1}\ldots\mathbf{x}_\tau]^T\in\mathbb{R}^{W\times F},\quad y_\tau=P(\tau+H)
\end{equation}
with $F=5$ features, split chronologically 70/15/15 (train $\approx$73{,}560, val $\approx$15{,}766, test $\approx$15{,}766). Min-Max normalization is fit per-condition on the training partition; metrics use de-normalized (MW) values against pristine ground truth.

\noindent\textbf{Proposed Architecture.} 1D-CNN Stage (64 filters, $k=3$, $s=1$, same padding, BatchNorm+ReLU):
\begin{equation}
\mathbf{C}_\tau=\text{ReLU}(\text{BatchNorm}(\mathbf{W}_c*\mathbf{X}_\tau+\mathbf{b}_c))\in\mathbb{R}^{W\times64}
\end{equation}
LSTM Stage (2-layer stacked, $d_h=64$):
\begin{equation}
\mathbf{h}_t,\mathbf{c}_t=\text{LSTM}(\mathbf{C}_\tau(t),\mathbf{h}_{t-1},\mathbf{c}_{t-1})
\end{equation}
Dense Head (32-unit ReLU then linear output):
\begin{equation}
\hat{y}_\tau=\mathbf{w}_o^T(\text{ReLU}(\mathbf{W}_d\mathbf{h}_W+\mathbf{b}_d))+b_o
\end{equation}
Identical architecture/hyperparameters are used across all conditions and for the plain LSTM baseline (CNN stage omitted), so differences are attributable to input condition and/or CNN presence.

\noindent\textbf{Training and Statistical Protocol.} Training uses Adam ($\eta=2\times10^{-3}$, weight decay $10^{-5}$), batch $B=256$, MSE loss, 12 epochs with validation-loss checkpointing, repeated across \textbf{5 seeds} $\{42,123,7,2024,99\}$ (datasets fixed, only training stochasticity varies). Persistence and AutoReg are deterministic. Three paired comparisons (Proposed vs.\ Plain LSTM, vs.\ Pristine, vs.\ Oracle-Cleaned) are tested via paired $t$-test across 5 seeds, with \textbf{Holm--Bonferroni correction} controlling the familywise error rate. Since a non-significant difference alone cannot establish equivalence, \textbf{TOST} tests whether the Proposed-vs.-Pristine difference falls within a pre-specified $\pm0.10$~MW margin ($\approx0.14\%$ of plant capacity).

\section{Experimental Results and Discussion}
\subsection{Benchmark Results}
Performance is evaluated on the de-normalized test set ($N_{\text{test}}\approx15{,}766$) via RMSE, MAE, $R^2$, and \textbf{Signal Recovery}, the fraction of corruption-induced error closed relative to pristine:
\begin{equation}
\text{Signal Recovery}=\frac{\text{RMSE}_{\text{raw}}-\text{RMSE}_{\text{cond}}}{\text{RMSE}_{\text{raw}}-\text{RMSE}_{\text{pristine}}}\times100\%
\label{eq:sr}
\end{equation}
Since the raw-to-pristine denominator is small ($\approx$0.31~MW), we treat Signal Recovery as illustrative only, relying on the paired $t$-tests/TOST below -- not the percentage -- for equivalence claims. Table~\ref{tab:results} summarizes multi-seed performance across all seven conditions.

\begin{table}[t]
\caption{Multi-seed ($n=5$) performance benchmark on the NREL 70~MW test set ($W{=}24$, $H{=}6$, 30-min-ahead). Deep rows: mean $\pm$ std.\ dev.\ over seeds $\{42,123,7,2024,99\}$; Persistence/AutoReg are deterministic. $^\dagger$Signal Recovery is a point estimate (Eq.~\ref{eq:sr}). $^*$Supported by TOST at $\pm0.10$~MW ($p\approx0.05$), not the percentage alone. Proposed-vs.-Plain-LSTM: raw $t=-3.89$, $p=0.018$; Holm $p=0.053$.}
\label{tab:results}
\centering
\footnotesize
\begin{tabular}{@{}lccccc@{}}
\toprule
\textbf{Model / Condition} & \textbf{RMSE (MW)} & \textbf{MAE (MW)} & \textbf{$R^2$} & \textbf{\%Impr.\ vs.\ Raw} & \textbf{Sig.\ Rec.$^\dagger$} \\
\midrule
CNN-LSTM (Pristine, ceiling)   & $3.58\pm0.06$ & $1.67\pm0.11$ & $0.949\pm0.002$ & +7.9\% & 100.0\% (ref.) \\
CNN-LSTM (Oracle-Cleaned)      & $3.58\pm0.05$ & $1.67\pm0.11$ & $0.949\pm0.002$ & +7.9\% & 100.0\%$^*$ \\
\textbf{CNN-LSTM (Proposed)}   & $\mathbf{3.57\pm0.04}$ & $\mathbf{1.64\pm0.08}$ & $\mathbf{0.949\pm0.001}$ & \textbf{+8.2\%} & \textbf{100.0\%}$^*$ \\
Plain LSTM (IF-Cleaned)        & $3.77\pm0.11$ & $1.78\pm0.12$ & $0.943\pm0.003$ & +3.0\% & 38.3\% \\
AutoReg (diurnal lag)          & 4.83 & 2.76 & 0.907 & $-$24.3\% & 0.0\% \\
Persistence (naive)            & 4.93 & 2.35 & 0.903 & $-$26.7\% & 0.0\% \\
CNN-LSTM (Raw, ablation)       & $3.89\pm0.11$ & $1.97\pm0.17$ & $0.940\pm0.003$ & Baseline & 0.0\% (ref.) \\
\bottomrule
\end{tabular}
\end{table}

\begin{figure}[t]
\centering
\begin{subfigure}{0.49\textwidth}
\includegraphics[width=\linewidth]{figure1_predicted_vs_actual.png}
\caption{Forecast vs.\ ground truth, 24-h test segment (seed 42).}
\end{subfigure}\hfill
\begin{subfigure}{0.49\textwidth}
\includegraphics[width=\linewidth]{figure2_rmse_bar_chart.png}
\caption{Error metrics across all seven conditions ($\pm1$ std.\ dev.).}
\end{subfigure}
\caption{Forecast tracking and aggregate error hierarchy. Pristine, oracle, and proposed trajectories closely overlap; raw-corrupted deviates visibly during ramps. See Table~\ref{tab:stats} for significance testing.}
\label{fig:pva_bar}
\end{figure}

\subsection{Analysis and Discussion}
\noindent\textbf{Impact of anomaly-aware preprocessing.} The four-condition ablation isolates preprocessing from architecture under identical hyperparameters. Raw-corrupted CNN-LSTM achieves RMSE $=3.89\pm0.11$~MW ($R^2=0.940$) versus a pristine ceiling of $3.58\pm0.06$~MW ($R^2=0.949$); the proposed IF-cleaned model reduces this to $3.57\pm0.04$~MW (8.2\% mean reduction). The gap is modest because at $\sim$2\% corruption and a 30-min horizon, the network's 2-hour context window already partially compensates for isolated corrupted points. IF-cleaned and oracle-labeled conditions are nearly identical ($3.57\pm0.04$ vs.\ $3.58\pm0.05$~MW), suggesting imperfect detection does not materially bottleneck the pipeline at this corruption rate -- interpolation quality, not detection accuracy, appears to govern recoverable performance, corroborated by the non-significant paired test below.

\noindent\textbf{Statistical validation.} Table~\ref{tab:stats} reports the three primary paired comparisons.

\begin{table}[t]
\caption{Paired $t$-tests for the proposed CNN-LSTM, Holm--Bonferroni corrected.}
\label{tab:stats}
\centering
\footnotesize
\begin{tabular}{@{}lcccc@{}}
\toprule
\textbf{Comparison} & \textbf{$\Delta$ (MW)} & \textbf{$t$} & \textbf{Raw $p$} & \textbf{Holm $p$} \\
\midrule
vs.\ Plain LSTM     & $-0.2015$ & $-3.892$ & $0.0177$ & $0.0530$ \\
vs.\ Pristine       & $-0.0122$ & $-0.294$ & $0.7837$ & $1.000$ \\
vs.\ Oracle-Cleaned & $-0.0108$ & $-0.280$ & $0.7934$ & $1.000$ \\
\bottomrule
\end{tabular}
\end{table}

The proposed model achieved lower RMSE than plain LSTM in every one of 5 paired seeds ($-0.148,-0.204,-0.091,-0.395,-0.170$~MW; raw $p=0.018$). After Holm correction, $p=0.053$ no longer clears 0.05; we report this as \emph{directionally consistent but not conclusively significant}. Comparisons against Pristine and Oracle-Cleaned remain far from significant (Holm $p=1.000$ both).

\noindent\textbf{TOST equivalence.} At a $\pm0.10$~MW margin: $\bar{d}=-0.0122$~MW ($s_d=0.0928$, $SE=0.0415$~MW); $t_1=+2.116$, $p_1=0.0509$; $t_2=-2.703$, $p_2=0.0270$; TOST $p=0.0509$. The 90\% CI is $[-0.1007,+0.0763]$~MW (equivalence confirmed at $\pm0.101$~MW, $p=0.049$). We describe the proposed pipeline as \emph{statistically equivalent to the pristine ceiling within an approximately $\pm0.10$~MW margin}, not unqualified 100\% recovery. Notably, the proposed CNN-LSTM shows the lowest RMSE std.\ dev.\ of any deep model ($\pm0.0393$~MW vs.\ $\pm0.1098$ plain LSTM, $\pm0.1097$ raw-data), suggesting cleaning also stabilizes optimization across seeds.

\noindent\textbf{Classical baselines.} At 30 min ahead, persistence and diurnal-lag AutoReg ($t{-}1,t{-}2,t{-}288$) degrade substantially versus deep approaches (4.93, 4.83~MW; $\sim$27\%, $\sim$26\% RMSE gaps vs.\ proposed), a larger margin than at a single-step horizon, since the near-future-resembles-recent-past assumption breaks down at 30~min.

\begin{figure}[t]
\centering
\begin{subfigure}{0.49\textwidth}
\includegraphics[width=\linewidth]{figure3_error_distribution.png}
\caption{Residual error distribution ($\hat y - y$, seed 42).}
\end{subfigure}\hfill
\begin{subfigure}{0.49\textwidth}
\includegraphics[width=\linewidth]{figure4_anomaly_detection_demo.png}
\caption{IF detection (red `X's) and reconstruction, 24-h window.}
\end{subfigure}
\caption{Residual error tails (proposed tracks pristine closely; raw shows wider tails) and a representative detection/reconstruction case (full-test detector: 72.21\% recall, 57.48\% precision; not a guarantee of flawless detection).}
\label{fig:err_demo}
\end{figure}

\section{Conclusion and Future Work}
We presented an anomaly-aware short-term solar forecasting framework combining unsupervised IF cleaning with a hybrid 1D CNN-LSTM at a 30-min horizon, validated via paired $t$-tests, Holm--Bonferroni correction, and TOST equivalence testing. Decomposing evaluation into pristine, raw-corrupted, IF-cleaned, and oracle-cleaned conditions, the proposed pipeline reaches a Signal Recovery point estimate of 100.0\%, formally supported by TOST showing statistical indistinguishability from the pristine ceiling within $\approx\pm0.10$~MW ($p\approx0.05$); detection quality does not measurably bottleneck this recovery relative to oracle knowledge (Holm-adjusted $p=1.000$). Across 5 seeds the proposed CNN-LSTM achieves RMSE $=3.57\pm0.04$~MW, beating plain LSTM in every paired seed, though after Holm correction this advantage does not reach conventional significance ($p=0.053$). Both deep models substantially and significantly outperform persistence and AutoReg baselines; at this horizon, preprocessing quality is the more strongly evidenced contributor to accuracy.

\noindent\textbf{Limitations and future work.} (1) $n=5$ seeds underpowers conclusive resolution of the borderline Plain-LSTM comparison (Holm $p=0.053$) and tightening the TOST margin below $\pm0.10$~MW; $n=20$--30 is needed. (2) Evaluation uses one 70~MW UPV facility; extension across climates and DPV fleets is needed. (3) Extending $H=6$ to day-ahead ($H=288$) to test whether the CNN-LSTM's advantage persists. (4) Bidirectional LSTMs or diffusion-based alternatives to linear interpolation for prolonged multi-hour outages.

\begin{thebibliography}{14}

\bibitem{ref1} Hochreiter, S., Schmidhuber, J.: Long short-term memory. Neural Computation \textbf{9}(8), 1735--1780 (1997)

\bibitem{ref2} Mellit, A., Pavan, A.M., Lughi, V.: Deep learning neural networks for short-term photovoltaic power forecasting. Renewable Energy \textbf{172}, 276--288 (2021)

\bibitem{ref3} Sobri, S., Koohi-Kamali, S., Rahim, N.A.: Solar photovoltaic generation forecasting methods: A review. Energy Conversion and Management \textbf{156}, 459--497 (2018)

\bibitem{ref4} Zang, H., et al.: Short-term solar power prediction based on CNN-LSTM hybrid neural network. Energy Reports \textbf{6}, 217--222 (2020)

\bibitem{ref5} Obiora, C.N., Ali, A., Hasan, A.N.: Hourly photovoltaic power forecasting using CNN-LSTM hybrid model. In: Proc. IEEE PES/IAS PowerAfrica Conf. (2021)

\bibitem{ref6} Qu, J., Qian, Z., Pei, Y.: Day-ahead hourly photovoltaic power forecasting using attention-based CNN-LSTM neural network. Energy \textbf{232}, 120996 (2021)

\bibitem{ref7} Boucetta, L.N., Amrane, Y., Chouder, A., Arezki, S., Kichou, S.: Enhanced forecasting accuracy of a grid-connected photovoltaic power plant: hybrid VMD and CNN-LSTM model. Energies \textbf{17}(7), 1781 (2024)

\bibitem{ref8} Zhang, J.E., et al.: Time series anomaly detection for smart grids: A survey. In: Proc. IEEE ISGT (2021)

\bibitem{ref9} Zhang, J., et al.: Baseline error analysis for solar irradiance and power forecasting. Solar Energy \textbf{122}, 109--124 (2015)

\bibitem{ref10} Abumohsen, M., Owda, A.Y., Owda, M., Abumihsan, A.: Hybrid ML model combining CNN-LSTM-RF for time series forecasting of solar power generation. e-Prime \textbf{9}, 100636 (2024)

\bibitem{ref11} Zhou, N., Shang, B., Xu, M., Peng, L., Feng, G.: Enhancing photovoltaic power prediction using a CNN-LSTM-attention hybrid model with Bayesian hyperparameter optimization. Global Energy Interconnection \textbf{7}(5), 667--681 (2024)

\bibitem{ref12} Liu, F.T., Ting, K.M., Zhou, Z.H.: Isolation forest. In: Proc. IEEE ICDM, pp. 413--422 (2008)

\bibitem{ref13} Jorry, V., et al.: An unsupervised kernel norm monitoring for fault detection in a time series photovoltaic system. arXiv:2607.01541 (2026)

% \bibitem{ref14} National Renewable Energy Laboratory (NREL): Solar Power Data for Integration Studies. Tech. Rep. NREL/TP-5500-65885 (2016)

\end{thebibliography}

\end{document}
